\section{Introduction}\label{sec:intro}

When a compact curve separates into several components, the scalar
Laplacian acquires small eigenvalues associated with functions that are
nearly constant on each component. For a bundle with parallel
endomorphisms, the same phenomenon occurs in the linearization of
Hermitian--Yang--Mills metric flow. We ask whether the evolving component
metrics admit a finite-dimensional description with an error tending to
zero as the curve degenerates.

We prove such a description for a trivial holomorphic bundle. Its finite
variables are positive Hermitian matrices, coupled across each node by
their Riemannian logarithm. Component areas supply the masses. The
equation is nonlinear even for the trivial bundle: metrics on three
components need not admit a simultaneous diagonalization. The comparison
holds for actual HYM solutions on intervals whose physical length grows
like the logarithm of the inverse smoothing parameter.

\subsection{The family and finite equation}
Let $\pi:\cX\to\Delta$ be a proper holomorphic map from a smooth complex
surface to a disc, with connected fibres. Assume that $\cX$ carries a
smooth K\"ahler form $\omega$, that $X_t=\pi^{-1}(t)$ is smooth for
$t\ne0$, and that $X_0=C_1\cup\cdots\cup C_N$ is reduced, with smooth
irreducible components and only ordinary nodes. Each node joins two
distinct components. Let $\mathcal E$ be the set of nodes, regarded as
edges of the component graph, retaining repeated edges. Put
\[
 A_i=\int_{C_i}\omega>0,\qquad L=L(t)=\log(1/|t|).
\]
All statements are uniform in the argument of sufficiently small $t$.
The induced metric on $X_t$ stays fixed during each flow.

Fix the trivial holomorphic rank-$r$ bundle, $r\ge2$, with its global
frame and determinant metric one. Write
\[
 \cP_r=\{H=H^*>0:\det H=1\}.
\]
Its symmetric-space metric and distance are
\begin{equation}\label{eq:target-metric}
 g_H(U,V)=\tr(H^{-1}UH^{-1}V),\qquad
 d(H,K)^2=\tr\bigl(\log(H^{-1/2}KH^{-1/2})^2\bigr).
\end{equation}
The principal matrix logarithm $\log(H^{-1}K)$ is self-adjoint for $H$;
it need not be Hermitian in the fixed frame.

Given $\mathbf H^0\in\cP_r^N$, define $\mathbf H(\tau)$ by
\begin{equation}\label{eq:finite-flow}
 A_iH_i^{-1}\frac{dH_i}{d\tau}
 =2\pi\sum_{p:\,i\leftrightarrow j}\log(H_i^{-1}H_j),
 \qquad H_i(0)=H_i^0.
\end{equation}
There is one summand for each incident node. This is the negative
gradient flow of
\begin{equation}\label{eq:graph-energy}
 \mathscr E(\mathbf H)=\pi\sum_{p:\,i\leftrightarrow j}d(H_i,H_j)^2
\end{equation}
for the product metric with weights $A_i$, each edge occurring once in
the energy. It has a global solution in a compact subset of $\cP_r^N$
(Lemma~\ref{lem:finite-flow}). For later comparison put
\begin{equation}\label{eq:sigma}
 \sigma(H,K)=\tr(H^{-1}K)+\tr(K^{-1}H)-2r\ge d(H,K)^2.
\end{equation}

\subsection{Actual-flow tracking}
Our HYM convention is
\begin{equation}\label{eq:hym-intro}
 H^{-1}\partial_sH=-2i\Lambda_{\omega_t}F_H,
 \qquad F_H=\bar\partial(H^{-1}\partial H).
\end{equation}
The degree and central curvature constant are zero. The variable $s$
is physical time and $\tau=s/L$ is slow time. For compact
$K_i\Subset C_i\setminus\operatorname{Sing}X_0$, let $K_i(t)\subset X_t$
be its image under the identifications of Lemma~\ref{lem:plumbing}.
All global integrals use $dA_t=\omega|_{X_t}$.

\begin{theorem}\label{thm:main}
For every initial tuple $\mathbf H^0$ there are smooth determinant-one
metrics $\widehat H_t(s)$ on the trivial bundle with the following
properties. For every fixed $T<\infty$, they and their inverses are
uniformly bounded for $0\le s\le TL$, and
\begin{equation}\label{eq:approx-thick-intro}
 \sup_{0\le s\le TL}\sup_{K_i(t)}
 d(\widehat H_t(s),H_i(s/L))\le C_{K_i,T}/L.
\end{equation}
Let $H_t^{\mathrm{in}}$ be smooth determinant-one metrics satisfying
$cI\le H_t^{\mathrm{in}}\le CI$ with constants independent of $t$, and
assume that
\[
 \eta_t=\int_{X_t}\sigma(H_t^{\mathrm{in}},\widehat H_t(0))\,dA_t
 \longrightarrow0.
\]
The actual HYM solution $H_t(s)$ exists smoothly for all $s\ge0$, and
\begin{equation}\label{eq:global-main}
 \sup_{0\le s\le TL}\int_{X_t}d(H_t(s),\widehat H_t(s))^2\,dA_t
 \le\eta_t+C_T\frac{(1+\log L)^2}{L}.
\end{equation}
For each fixed compact $K_i$ one also has
\begin{equation}\label{eq:local-main}
 \sup_{1\le s\le TL}\sup_{K_i(t)}d(H_t(s),H_i(s/L))^2
 \le C_{K_i,T}\left(\eta_t+\frac{(1+\log L)^2}{L}\right).
\end{equation}
The constants may depend on the fixed family, rank, initial tuple,
$c,C$ and $T$, but not on $t$ or the spatial derivatives of the initial
metrics. The interval in \eqref{eq:local-main} is nonempty for small $t$.
\end{theorem}

The comparison metrics are constructed from component Green functions,
matrix exponentials and geodesics. In particular the initial class is
nonempty: take $H_t^{\mathrm{in}}=\widehat H_t(0)$. Uniform bounds and
$L^2$ convergence of the initial metrics to $H_i^0$ on every fixed
compact component subset also imply $\eta_t\to0$, by the node-area bound.
Initial derivatives and bounded initial neck values are unrestricted.

The global estimate has no waiting time. The pointwise estimate has a
fixed physical-time transient, equal to $1/L$ in slow time. It concerns
fixed compact subsets away from the nodes and makes no assertion of
uniform pointwise tracking through the shrinking neck. The theorem also
does not assert tracking uniform as $T\to\infty$.

\subsection{Linearized evolution and a compact example}
At the constant identity metric, \eqref{eq:hym-intro} linearizes to the
scalar heat equation on trace-free Hermitian matrix entries. Write $Q$
for the component graph Laplacian, $A=\operatorname{diag}(A_i)$ and
$B=2\pi A^{-1}Q$. The scalar theorem of \cite{wuebben-curves} gives
$N-1$ small positive eigenvalues and a complementary gap.
Theorem~\ref{thm:operator} identifies the low space by component averages
and proves convergence of the rescaled heat evolution to $e^{-\tau B}$
in operator norm. Tensoring gives $(N-1)(r^2-1)$ small positive modes
and $r^2-1$ exact parallel modes. The finite matrix flow has the same
linearization at a constant tuple.

The cubic family $xyz=t(x^3+y^3+z^3)$ has a central triangle with
component areas one in the normalized Fubini--Study metric. The tuple
\begin{equation}\label{eq:example-tuple-intro}
 I,\qquad\begin{pmatrix}2&0\\0&1/2\end{pmatrix},\qquad
 \begin{pmatrix}5/3&4/3\\4/3&5/3\end{pmatrix}
\end{equation}
cannot be simultaneously diagonalized by a fixed complex congruence.
Section~\ref{sec:example} verifies the geometry and the genuinely
noncommuting initial data. Its rank-two positive linearized cluster
has six modes, each asymptotic to $6\pi/L$.

\subsection{Predecessors and proof}
Donaldson's heat flow and metric comparison provide the analytic
framework \cite{donaldson,hkkp}. Haiden--Katzarkov--Kontsevich--Pandit
describe algebraically determined asymptotics on a fixed curve, and
Lam treats broader fixed compact K\"ahler backgrounds under his stated
hypotheses \cite{lam}. The local spectral and phase comparisons in
\cite{wuebben-local,wuebben-phase} also concern fixed smooth
backgrounds; their regularity constants cannot be assumed uniform
under a nodal degeneration.

Daskalopoulos--Wentworth already study actual Yang--Mills flow near
the boundary of Teichm\"uller space \cite{dw}. Their uniform long-time
argument uses a positive adjoint spectral gap when the limiting
representations on the normalization components are irreducible;
their spectral analysis also detects small adjoint eigenvalues.
The componentwise parallel endomorphisms of the trivial bundle lie
outside that gap regime. The present theorem computes their finite
nonlinear evolution on logarithmic time and reconstructs actual flows.

The intrinsic squared-distance graph energy and its logarithmic
gradient are established constructions \cite{tav}. The assertion here
is their occurrence with these geometric masses and capacities in an
HYM degeneration. The HYM equation contains a commutator term absent
from ordinary symmetric-space harmonic-map heat flow
\eqref{eq:real-hym}. We retain that term and use exact stationary
one-variable profiles only on the necks.

Harmonic-map degeneration has a separate compactness theory.
Zhu \cite{zhu} proves energy identities and compactness criteria for
bounded-energy sequences of harmonic maps from degenerating hyperbolic
surfaces to a compact target. Rupflin--Topping--Zhu \cite{rtz} study
Teichm\"uller harmonic map flow, in which both the map and its hyperbolic
domain metric evolve; for global solutions degenerating only at infinite
time, they obtain subsequential limits given by branched minimal
immersions or constant maps on compactified components. Here each induced
fibre metric remains fixed during its HYM flow, and the comparison is
uniform on bounded logarithmic time intervals. The commutator in
\eqref{eq:real-hym} also distinguishes the two evolution equations.

The proof adds Green-function corrections to the component metrics
and matches them to geodesic neck profiles. The resulting curvature
residual has space-time $L^1$ norm tending to zero on $0\le s\le TL$.
A forced Donaldson trace inequality proves global tracking, and local
scalar parabolic estimates give pointwise convergence away from the
nodes. The argument does not require a global unweighted Schauder
estimate uniform under degeneration. The separate spectral comparison
provides uniform Hilbert-space inverse and heat bounds on the complement.
